This report shows the effectiveness of RQL combustion in reducing \ce{NOx} emissions from ammonia-fueled systems. The emmision of \ce{NOx} is reduced by using a 3 phase reactor. The first phase is burning ammonia in a fuel-rich environment, then adding secondary air, and finally completing the combustion in a lean zone. The use of Cantera and Ansys validates the results. Cantera showed a significant drop in \ce{NOx} from the rich zone to the lean zone exiting with 55.47 ppm, while ansys simulations show higher outlet levels of \ce{NOx} (2058 ppm), confirming the general trend and highlighting the role of temperature and air mixing in the control of emission. These findings indicate that RQL combustion is a viable approach for enabling cleaner and more controlled ammonia combustion in low-emission gas turbine applications.